\begin{table}[!h]
\centering
\caption{\label{tab:gp_summary_combined}Summary of Gram Panchayats by Election Year}
\centering
\fontsize{8}{10}\selectfont
\begin{threeparttable}
\begin{tabular}[t]{lrrrr}
\toprule
 & 2005 & 2010 & 2015 & 2020/2021\\
\midrule
\addlinespace[0.3em]
\multicolumn{5}{l}{\textit{\textbf{Panel A: Rajasthan}}}\\
\hspace{1em}Identified GPs & 9,156 & 9,166 & 9,862 & 11,311\\
\hspace{1em}Blocks & 237 & 249 & 300 & 352\\
\hspace{1em}Districts & 32 & 33 & 33 & 33\\
\hspace{1em}Women reserved (\%) & 33.5 & 47.6 & 48.4 & 48.6\\
\addlinespace[0.3em]
\multicolumn{5}{l}{\textit{\textbf{Panel B: Uttar Pradesh}}}\\
\hspace{1em}Identified GPs & 51,404 & 51,757 & 58,994 & 49,772\\
\hspace{1em}Blocks & 842 & 843 & 823 & 728\\
\hspace{1em}Districts & 70 & 72 & 75 & 67\\
\hspace{1em}Women reserved (\%) & 46.6 & 33.1 & 33.8 & 33.7\\
\bottomrule
\end{tabular}
\begin{tablenotes}
\item \textit{Note: } 
\item Identified GPs are distinct complete district--block--GP names in the source files, using cleaned GP names. Missing identities are excluded from GP counts; repeated names within a block count once. Reservation shares use identified GPs with known, non-conflicting reservation. District and block counts use observed names. The last column contains Rajasthan 2020 and UP 2021. Coverage changes reflect both administrative changes and source availability.
\end{tablenotes}
\end{threeparttable}
\end{table}
